\chapter{Quadratic characters and the criterion}
\label{chap:characters}

\section{Quadratic characters of \texorpdfstring{$\QQ$}{Q}}

\begin{definition}[Quadratic characters]
  \label{def:qchar}
  \lean{MulChar.apply_eq_one_or_eq_neg_one, MulChar.isQuadratic_int,
    MulChar.apply_ne_neg_one_of_isSquare, MulChar.apply_inv, MulChar.apply_div, MulChar.apply_sq,
    MulChar.apply_mul_sq, MulChar.apply_div_sq}
  \leanok
  A \emph{quadratic character} of $\QQ$ is a homomorphism $\chi \colon \QQ^\times \to \{\pm 1\}$,
  extended to $\QQ$ by $\chi(0) = 0$. Equivalently, a multiplicative map $\chi \colon \QQ \to \ZZ$
  with $\chi(1) = 1$ and $\chi(0) = 0$; its values on $\QQ^\times$ are units of $\ZZ$, i.e.\
  $\pm 1$. Quadratic characters form a group under pointwise multiplication, and
  $\chi(x y^2) = \chi(x)$ for $y \ne 0$; in particular $\chi(x) \ne -1$ when $x$ is a square, and
  $\chi(1/x) = \chi(x)$.
  \formalnote{Mathlib's \lname{MulChar ℚ ℤ}, a \lname{CommGroup}; it corresponds to
  \lname{ℚˣ →* ℤˣ} by \lname{MulChar.equivToUnitHom}, i.e.\ to the characters of
  $\QQ^\times/\QQ^{\times 2}$. The general lemmas are stated for a character with values in $\ZZ$
  of a commutative monoid (\lname{MulChar.apply\_eq\_one\_or\_eq\_neg\_one} at units, so every
  such character is quadratic) or of a commutative group with zero (\lname{MulChar.apply\_inv},
  \lname{MulChar.apply\_div}, \lname{MulChar.apply\_mul\_sq}, \lname{MulChar.apply\_div\_sq}).
  \lname{MulChar.apply\_ne\_neg\_one\_of\_isSquare}, for values in an ordered ring, is the whole
  reason a character detects non-squares: $\chi(y^2) = \chi(y)^2 \ge 0$.}
\end{definition}

\begin{definition}[The characters used]
  \label{def:chars}
  \uses{def:qchar}
  \lean{Rat.mulCharOfInt, Rat.mulCharOfInt_apply, Rat.mulCharOfInt_intCast,
    Rat.map_num_mul_map_den_mul, Rat.signChar, Rat.signChar_apply, Rat.signChar_apply_of_pos,
    Rat.signChar_apply_of_neg, Rat.valChar, Rat.valChar_intCast, Rat.unitPartChar,
    Rat.unitPartChar_intCast, Rat.unitPartChar_natCast_self, ZMod.isUnit_intCast_of_not_dvd}
  \leanok
  For $x \in \QQ^\times$ and a prime $p$ let $x_p = p^{-v_p(x)} x$ be the prime-to-$p$ part.
  \begin{enumerate}
    \item $\sgn(x)$, the sign.
    \item $\val_p(x) = (-1)^{v_p(x)}$, for a prime $p$.
    \item For a prime $p$, a power $N$ of $p$ and a character $\psi$ of $(\ZZ/N\ZZ)^\times$ with
          values $\pm 1$ (extended by $0$ to $\ZZ/N\ZZ$): $\chi_\psi(x) = \psi(x_p \bmod N)$,
          where $x_p$ is reduced modulo $N$ as a fraction with denominator prime to $p$.
          Special cases:
          $\chi_p = \chi_\psi$ for an odd prime $p$, $N = p$ and $\psi$ the Legendre symbol;
          $\chi_4 = \chi_\psi$ for $N = 4$ and $\psi(-1) = -1$;
          $\chi_8 = \chi_\psi$ for $N = 8$ and $\psi(\pm 1) = 1$, $\psi(\pm 3) = -1$.
  \end{enumerate}
  Then $\chi_\psi(p) = 1$, and $\chi_\psi(m) = \psi(m \bmod N)$ for an integer $m$ prime to $p$;
  more generally $\chi_\psi(m/k) = \psi(m)\psi(k)$ for integers $m$, $k$ prime to $p$.
  We call $\val_p$, $\chi_p$ (resp.\ $\val_2$, $\chi_4$, $\chi_8$) and their products the
  \emph{characters at $p$}; they are induced by the characters of $\QQ_p^\times/\QQ_p^{\times 2}$.
  \formalnote{In the mirror file \lname{QuadraticIterates/Mathlib/NumberTheory/MulChar/Rat.lean}.
  All three are extensions $x \mapsto f(\mathrm{num}\,x) f(\mathrm{den}\,x)$ of multiplicative
  functions $f$ on $\ZZ$ that are $\pm 1$ at nonzero integers (\lname{Rat.mulCharOfInt}):
  \lname{Rat.signChar}, \lname{Rat.valChar p} and \lname{Rat.unitPartChar p ψ}, the last one
  defined for any \lname{ψ : MulChar (ZMod N) ℤ}, through $\psi(\operatorname{sign}(m) \cdot
  |m|/p^{v_p(m)})$, with value lemmas for $N = p^k$: \lname{Rat.unitPartChar\_intCast}
  ($p \nmid m$) and \lname{Rat.unitPartChar\_natCast\_self} ($\chi_\psi(p) = 1$). The value at
  $m/k$ follows from \lname{MulChar.apply\_div}. $\chi_p$, $\chi_4$, $\chi_8$ are
  \lname{Rat.unitPartChar} of \lname{quadraticChar (ZMod p)}, \lname{ZMod.χ₄} and \lname{ZMod.χ₈}.}
\end{definition}

\section{Patterns}

\begin{definition}[Patterns]
  \label{def:pattern}
  \uses{def:qchar}
  \lean{QuadraticIterates.HasCharPattern, QuadraticIterates.HasCharPattern.mul,
    QuadraticIterates.HasCharPattern.mono, QuadraticIterates.HasCharPattern.congr}
  \leanok
  Let $(v_m)_{m \ge 1}$ be a sequence of rational numbers, $n \ge 1$ and
  $\varepsilon \colon \ZZ_{\ge n} \to \{\pm 1\}$. The \emph{pattern $\varepsilon$ is realized
  from $n$} if some quadratic character $\chi$ satisfies $\chi(v_m) = \varepsilon_m$ for all
  $m \ge n$. Patterns realized from $n$ are closed under pointwise products (multiply the
  characters), contain the constant pattern $1$ if all $v_m$ are nonzero (the trivial
  character), and a pattern realized
  from $n$ is realized from every $n' \ge n$. We write patterns as strings of signs:
  $\chi(v_\bullet) = ({+}{-}{-}\cdots)$.
  \formalnote{\lname{HasCharPattern (v : ℕ → M) (n : ℕ) (ε : ℕ → ℤ) : Prop :=
  ∃ χ : MulChar M ℤ, ∀ m ≥ n, χ (v m) = ε m}, for a sequence in a commutative monoid $M$, with
  \lname{HasCharPattern.congr} for patterns agreeing from $n$ on. Values of $\varepsilon$ before
  $n$ are irrelevant. The constant pattern $1$ is not needed.}
\end{definition}

\begin{definition}[$\phi_n$, $\phi_0$ and $\tau$]
  \label{def:phin}
  \uses{def:pattern, def:cseq}
  \leanok
  For the sequence $(c_m)$ of $a \in \QQ$:
  \begin{itemize}
    \item $a$ \emph{has $\phi_n$} ($n \ge 1$) if the pattern $-1$ at $n$, $+1$ after $n$ is
          realized from $n$: a character $\phi_n$ with $\phi_n(c_n) = -1$ and
          $\phi_n(c_m) = 1$ for all $m > n$;
    \item $a$ \emph{has $\phi_0$} if the constant pattern $-1$ is realized from $1$:
          $\phi_0(c_m) = -1$ for all $m \ge 1$;
    \item $a$ \emph{has $\tau$} if the pattern $({+}{-}{-}\cdots)$ is realized from $1$:
          $\tau(c_1) = 1$ and $\tau(c_m) = -1$ for all $m \ge 2$.
  \end{itemize}
  \formalnote{The patterns are \lname{Pi.mulSingle n (-1)}, the constant \lname{-1} and
  \lname{-Pi.mulSingle 1 (-1)}; statements are written with \lname{HasCharPattern (cSeq a)}
  directly, without further definitions.}
\end{definition}

\section{The criterion}

\begin{proposition}[The criterion with characters]
  \label{prop:crit}
  \uses{def:pattern, def:twoindep}
  \lean{QuadraticIterates.twoIndependent_of_forall_hasCharPattern,
    QuadraticIterates.forall_nonempty_mulEquiv_of_forall_hasCharPattern}
  \leanok
  Let $(v_m)_{m \ge 1}$ be a sequence of rational numbers such that for every $n \ge 1$ the pattern
  $-1$ at $n$, $+1$ after $n$ is realized from $n$. Then $(v_m)_{m \ge 1}$ is $2$-independent.
  In particular, if $a$ has $\phi_n$ for every $n \ge 1$, then $(c_n)_{n \ge 1}$ is
  $2$-independent and $\rho_a$ is surjective.
  \formalnote{For a sequence in any commutative monoid, with conclusion
  \lname{TwoIndependent (fun m : ℕ ↦ v (m + 1))}.}
\end{proposition}
\begin{proof}
  \leanok
  \uses{thm:criterion-infinite}
  Let $S \subseteq \ZZ_{\ge 1}$ be finite and nonempty, $n = \min S$, and $\phi_n$ as in the
  hypothesis. Then $\phi_n(\prod_{k \in S} v_k) = \phi_n(v_n) \prod_{k \in S, k > n}
  \phi_n(v_k) = -1$, so the product is not a square (Definition~\ref{def:qchar}). The second
  statement follows with Theorem~\ref{thm:criterion-infinite}.
\end{proof}

\begin{lemma}[Characters from linear forms]
  \label{lem:dual}
  \uses{def:qchar, def:twoindep}
  \lean{QuadraticIterates.mulCharOfDual, QuadraticIterates.mulCharOfDual_apply}
  \leanok
  Let $f \colon \QQ^\times/\QQ^{\times 2} \to \FF_2$ be $\FF_2$-linear. Then
  $x \mapsto (-1)^{f([x])}$ on $\QQ^\times$, $0 \mapsto 0$, is a quadratic character.
  If $V$ is a subspace of $\QQ^\times/\QQ^{\times 2}$ and $[x] \notin V$, there is such an $f$
  with $f([x]) = 1$ and $f(V) = 0$.
  \formalnote{For any field $K$ in place of $\QQ$, with \lname{SquareClasses K} and
  \lname{sqClass}: \lname{mulCharOfDual (f : Module.Dual (ZMod 2) (SquareClasses K)) :
  MulChar K ℤ}; the separation is Mathlib's
  \lname{Submodule.exists\_dual\_map\_eq\_bot\_of\_notMem} (every module over a field is free).}
\end{lemma}
\begin{proof}
  \leanok
  The map is the composition of the class map, $f$, and the isomorphism
  $\FF_2 \cong \{\pm 1\}$; vector spaces over a field have enough linear forms to separate a
  vector from a subspace.
\end{proof}

\begin{proposition}[The converse]
  \label{prop:crit-converse}
  \uses{lem:dual}
  \lean{QuadraticIterates.hasCharPattern_of_twoIndependent}
  \leanok
  If $(v_m)_{m \ge 1}$ is $2$-independent, then for every $n \ge 1$ the pattern $-1$ at $n$,
  $+1$ after $n$ is realized from $n$.
  \formalnote{For a sequence in a field. The classes of $v_1, v_2, \dots$ are linearly
  independent (\lname{twoIndependent\_iff\_linearIndependent} holds for families indexed by any
  type), and \lname{LinearIndependent.notMem\_span\_image} gives $[v_n] \notin V$.}
\end{proposition}
\begin{proof}
  \leanok
  \uses{lem:dual, def:twoindep}
  All $v_m$ are nonzero. If $[v_n]$ lay in the span of the $[v_m]$, $m > n$, then $v_n$ times a
  finite product of later terms would be a square, contradicting $2$-independence. So
  Lemma~\ref{lem:dual} gives $f$ with $f([v_n]) = 1$ and $f([v_m]) = 0$ for $m > n$, and the
  character $(-1)^f$ realizes the pattern.
\end{proof}

\begin{theorem}[The criterion]
  \label{thm:tfae}
  \uses{thm:criterion-infinite, prop:crit, prop:crit-converse}
  For $a \in \QQ$ the following are equivalent:
  \begin{enumerate}
    \item $\Om_n \cong \Ct{n}$ for all $n$, i.e.\ $\rho_a$ is surjective;
    \item $(c_n)_{n \ge 1}$ is $2$-independent;
    \item $a$ has $\phi_n$ for every $n \ge 1$.
  \end{enumerate}
  \lean{QuadraticIterates.tfae_forall_nonempty_mulEquiv}
  \leanok
\end{theorem}
\begin{proof}
  \leanok
  \uses{thm:criterion-infinite, prop:crit, prop:crit-converse}
  (1)$\iff$(2) is Theorem~\ref{thm:criterion-infinite}, (3)$\Rightarrow$(2) is
  Proposition~\ref{prop:crit}, (2)$\Rightarrow$(3) is Proposition~\ref{prop:crit-converse}.
\end{proof}

\section{Combining characters}

The lemmas of this section produce $\phi_n$ from simpler patterns by multiplying characters.

\begin{lemma}[Two of three]
  \label{lem:klein}
  \uses{def:phin}
  \lean{QuadraticIterates.hasCharPattern_neg_mulSingle_of_neg_of_mulSingle,
    QuadraticIterates.hasCharPattern_mulSingle_of_neg_of_neg_mulSingle,
    QuadraticIterates.hasCharPattern_neg_of_mulSingle_of_neg_mulSingle}
  \leanok
  If $a$ has two of $\phi_0$, $\phi_1$, $\tau$, it has the third.
  \formalnote{The patterns $({-}{-}{-}\cdots)$, $({-}{+}{+}\cdots)$, $({+}{-}{-}\cdots)$ are the
  non-trivial elements of a Klein four-group under pointwise multiplication; three one-line
  consequences of \lname{HasCharPattern.mul}, stated for \lname{Pi.mulSingle k (-1)} with any
  $k$ and from any $n$.}
\end{lemma}
\begin{proof}
  \leanok
  The product of the patterns of any two of $\phi_0$, $\phi_1$, $\tau$ is the pattern of the
  third.
\end{proof}

\begin{lemma}[A jump after $n$]
  \label{lem:jump}
  \uses{def:phin}
  \lean{QuadraticIterates.hasCharPattern_mulSingle_of_neg_of_forall_lt}
  \leanok
  Let $n \ge 1$. If $a$ has $\phi_0$ and some character $\psi$ satisfies $\psi(c_n) \ne 0$ and
  $\psi(c_m) = -\psi(c_n)$ for all $m > n$, then $a$ has $\phi_n$.
\end{lemma}
\begin{proof}
  \leanok
  Take $\phi_n = \psi$ if $\psi(c_n) = -1$ and $\phi_n = \phi_0 \psi$ if $\psi(c_n) = 1$.
\end{proof}

\begin{lemma}[A period-two jump]
  \label{lem:jump2}
  \uses{def:phin}
  \lean{QuadraticIterates.hasCharPattern_mulSingle_of_neg_of_alternating}
  \leanok
  Let $n \ge 1$. Assume that $a$ has $\phi_0$, that some character $\psi$ satisfies
  $\psi(c_n) \ne 0$ and $\psi(c_{m+1}) = -\psi(c_m)$ for all $m \ge n$, and that some character
  $\chi$ has the pattern $(x, y, -x, y, -x, y, \dots)$ from $n$, i.e.\ $\chi(c_n) = x$,
  $\chi(c_{n+2k}) = -x$ for $k \ge 1$ and $\chi(c_{n+2k+1}) = y$ for $k \ge 0$, with
  $x, y \in \{\pm 1\}$. Then $a$ has $\phi_n$.
  \formalnote{The pattern of $\chi$ is stated as $\chi(c_{n+2}) = -\chi(c_n) \ne 0$,
  $\chi(c_{n+1}) \ne 0$, and $\chi(c_{m+2}) = \chi(c_m)$ for $m > n$.}
\end{lemma}
\begin{proof}
  \leanok
  \uses{lem:jump}
  If $y = -x$, then $\chi$ is $x$ at $c_n$ and $-x$ at all later terms. If $y = x$, then with
  $s = \psi(c_n)$ the character $\chi\psi$ has the pattern $(xs, -xs, -xs, \dots)$ from $n$.
  In both cases Lemma~\ref{lem:jump} applies.
\end{proof}
